\section{Research synthesis and implications for TIDE}
\subsection{Common questions across substrates}
The three substrate classes expose different parts of the same constructive
problem. Soft systems make causality and heredity inspectable; hardware makes
contact, manufacturing, and control errors unavoidable; wet systems expose
whether maintenance, compartmentalization, and information copying can operate
together chemically. No single demonstration should substitute for the
integration problem addressed by another class.

Across them, five questions remain useful. First, what is the reproducing
unit, and can that unit change? Second, which resources and environmental
services are supplied? Third, how does inherited organization affect future
opportunities? Fourth, which processes are controlled by an external optimizer
or experimenter? Finally, which observation would falsify the proposed
mechanism? These questions organize comparisons without requiring all systems
to use the same substrate or pursue the same objective.

For soft artificial life, a productive priority is to make mechanistic claims
sharper rather than to add more life-like vocabulary. New instructions,
environmental gradients, ecological rewards, or learned evaluators enlarge
the design space, but also add assumptions. Their value is established by
controlled contrasts that explain what becomes possible and what changes in
measured outcomes. Negative results can identify a missing expressive
capability, an unfavorable ecological regime, or a limitation in the observer.

\subsection{TIDE's present position}
TIDE currently has research documentation and an earlier local prototype;
the fresh C++ kernel has not been implemented. The project has selected an
individual-CPU local address-based VM, pointer-directed byte copying,
allocation that occupies public lattice space before division, and a two-phase
synchronous default. It begins with a supplied ancestor. Ecological labels
must not determine physical behavior. These are project decisions, not results
proved by the reviewed literature.

The instruction-level contract remains open. In particular, read visibility,
conflict resolution, newborn scheduling, reclamation, and mutation require
explicit transitions including failure effects. Working defaults such as
readable residual bytes and CPU-only costs remain experimental variables.
Choosing a reference CPU implementation and deterministic replay helps
separate implementation discrepancies from intended changes in the model.

\subsection{Candidate directions and discriminating tests}
The current design discussion asks whether channel-dependent boundaries can
be constructed and maintained by local dynamics. Candidate ``option B'' keeps
execution threads and per-site ownership as supplied primitives, while leaving
body extent, shape, territorial division, and multi-thread grouping open to
evolution. It therefore leaves some units given while seeking emergence at
other scales. It is a discussion direction, not an amendment to the agreed
kernel contract. The second literature pass supplies design constraints and
failure cases without assessing novelty.

One useful controlled comparison from the earlier pass remains:
\begin{quote}
Starting from a supplied ancestor, how do local data-reading and
instruction-fetch permissions affect the appearance and persistence of
functionally distinct replication strategies?
\end{quote}
This is a proposed research question. Self-only, neighbor-data access, and neighbor-data-plus-fetch
conditions could be compared while holding write permissions and other rules
fixed. The ancestor must be viable across the relevant comparison. Mutation-free
examples establish expressive capacity before evolutionary runs begin.

Extending that comparison to constructed boundaries requires additional
decisions: what changes each channel's transmission, which material or work
is consumed, how decay and repair affect access, and how simultaneous requests
are resolved. Group formation then needs separate join, leave, fission,
death, and transfer semantics. An observer's inferred individuality should
remain an analysis result rather than an input granting physical privileges.

Two subordinate questions are worth retaining. Does readable residual code
support heritable replication opportunities after its producer dies, or merely
alter occupation and transient copying? Are observed dependencies robust to
the specified sequential-update comparison and alternative conflict policies?
The first requires causal removal/replacement of recorded bytes; the second
requires matched execution budgets and distributional comparisons across
seeds. Neither has been answered by TIDE experiments.

The minimum evidence package should identify all seeds and configurations,
archive the kernel revision, retain lifecycle and material-provenance records,
and report stagnation and extinction alongside surviving trajectories.
Implementation acceptance concerns semantics, locality, accounting, and
replay. Research acceptance must also allow a well-supported negative answer;
observing a desired ecological role should not be an engineered requirement.

\section{Conclusion}
Artificial life contains a family of constructive experiments rather than a
single route from simple rules to complete life. Self-organization,
self-reproduction, heredity, adaptation, and continuing innovation remain
distinct achievements. Digital programs and automata are especially useful
for tracing how inherited organization changes its own reproductive pathways;
robotic and biochemical work clarifies which constraints disappear when that
organization is simulated. The immediate opportunity for TIDE is a precise,
replayable investigation of replication access and ecological feedback, with
measurements that distinguish changing bytes from changing strategies. The
literature supports this as a tractable research direction while leaving its
outcome open.
